\IfFileExists{stacks-project.cls}{%
\documentclass{stacks-project}
}{%
\documentclass{amsart}
}

% For dealing with references we use the comment environment
\usepackage{verbatim}
\newenvironment{reference}{\comment}{\endcomment}
%\newenvironment{reference}{}{}
\newenvironment{slogan}{\comment}{\endcomment}
\newenvironment{history}{\comment}{\endcomment}

% For commutative diagrams we use Xy-pic
\usepackage[all]{xy}

% We use 2cell for 2-commutative diagrams.
\xyoption{2cell}
\UseAllTwocells

% We use multicol for the list of chapters between chapters
\usepackage{multicol}

% This is generally recommended for better output
\usepackage{lmodern}
\usepackage[T1]{fontenc}

% For cross-file-references
\usepackage{xr-hyper}

% Package for hypertext links:
\usepackage{hyperref}

% For any local file, say "hello.tex" you want to link to please
% use \externaldocument[hello-]{hello}
\externaldocument[introduction-]{introduction}
\externaldocument[conventions-]{conventions}
\externaldocument[sets-]{sets}
\externaldocument[categories-]{categories}
\externaldocument[topology-]{topology}
\externaldocument[sheaves-]{sheaves}
\externaldocument[sites-]{sites}
\externaldocument[stacks-]{stacks}
\externaldocument[fields-]{fields}
\externaldocument[algebra-]{algebra}
\externaldocument[brauer-]{brauer}
\externaldocument[homology-]{homology}
\externaldocument[derived-]{derived}
\externaldocument[simplicial-]{simplicial}
\externaldocument[more-algebra-]{more-algebra}
\externaldocument[smoothing-]{smoothing}
\externaldocument[modules-]{modules}
\externaldocument[sites-modules-]{sites-modules}
\externaldocument[injectives-]{injectives}
\externaldocument[cohomology-]{cohomology}
\externaldocument[sites-cohomology-]{sites-cohomology}
\externaldocument[dga-]{dga}
\externaldocument[dpa-]{dpa}
\externaldocument[sdga-]{sdga}
\externaldocument[hypercovering-]{hypercovering}
\externaldocument[schemes-]{schemes}
\externaldocument[constructions-]{constructions}
\externaldocument[properties-]{properties}
\externaldocument[morphisms-]{morphisms}
\externaldocument[coherent-]{coherent}
\externaldocument[divisors-]{divisors}
\externaldocument[limits-]{limits}
\externaldocument[varieties-]{varieties}
\externaldocument[topologies-]{topologies}
\externaldocument[descent-]{descent}
\externaldocument[perfect-]{perfect}
\externaldocument[more-morphisms-]{more-morphisms}
\externaldocument[flat-]{flat}
\externaldocument[groupoids-]{groupoids}
\externaldocument[more-groupoids-]{more-groupoids}
\externaldocument[etale-]{etale}
\externaldocument[chow-]{chow}
\externaldocument[intersection-]{intersection}
\externaldocument[pic-]{pic}
\externaldocument[weil-]{weil}
\externaldocument[adequate-]{adequate}
\externaldocument[dualizing-]{dualizing}
\externaldocument[duality-]{duality}
\externaldocument[discriminant-]{discriminant}
\externaldocument[derham-]{derham}
\externaldocument[local-cohomology-]{local-cohomology}
\externaldocument[algebraization-]{algebraization}
\externaldocument[curves-]{curves}
\externaldocument[resolve-]{resolve}
\externaldocument[models-]{models}
\externaldocument[functors-]{functors}
\externaldocument[equiv-]{equiv}
\externaldocument[pione-]{pione}
\externaldocument[etale-cohomology-]{etale-cohomology}
\externaldocument[proetale-]{proetale}
\externaldocument[relative-cycles-]{relative-cycles}
\externaldocument[more-etale-]{more-etale}
\externaldocument[trace-]{trace}
\externaldocument[crystalline-]{crystalline}
\externaldocument[spaces-]{spaces}
\externaldocument[spaces-properties-]{spaces-properties}
\externaldocument[spaces-morphisms-]{spaces-morphisms}
\externaldocument[decent-spaces-]{decent-spaces}
\externaldocument[spaces-cohomology-]{spaces-cohomology}
\externaldocument[spaces-limits-]{spaces-limits}
\externaldocument[spaces-divisors-]{spaces-divisors}
\externaldocument[spaces-over-fields-]{spaces-over-fields}
\externaldocument[spaces-topologies-]{spaces-topologies}
\externaldocument[spaces-descent-]{spaces-descent}
\externaldocument[spaces-perfect-]{spaces-perfect}
\externaldocument[spaces-more-morphisms-]{spaces-more-morphisms}
\externaldocument[spaces-flat-]{spaces-flat}
\externaldocument[spaces-groupoids-]{spaces-groupoids}
\externaldocument[spaces-more-groupoids-]{spaces-more-groupoids}
\externaldocument[bootstrap-]{bootstrap}
\externaldocument[spaces-pushouts-]{spaces-pushouts}
\externaldocument[spaces-chow-]{spaces-chow}
\externaldocument[groupoids-quotients-]{groupoids-quotients}
\externaldocument[spaces-more-cohomology-]{spaces-more-cohomology}
\externaldocument[spaces-simplicial-]{spaces-simplicial}
\externaldocument[spaces-duality-]{spaces-duality}
\externaldocument[formal-spaces-]{formal-spaces}
\externaldocument[restricted-]{restricted}
\externaldocument[spaces-resolve-]{spaces-resolve}
\externaldocument[formal-defos-]{formal-defos}
\externaldocument[defos-]{defos}
\externaldocument[cotangent-]{cotangent}
\externaldocument[examples-defos-]{examples-defos}
\externaldocument[algebraic-]{algebraic}
\externaldocument[examples-stacks-]{examples-stacks}
\externaldocument[stacks-sheaves-]{stacks-sheaves}
\externaldocument[criteria-]{criteria}
\externaldocument[artin-]{artin}
\externaldocument[quot-]{quot}
\externaldocument[stacks-properties-]{stacks-properties}
\externaldocument[stacks-morphisms-]{stacks-morphisms}
\externaldocument[stacks-limits-]{stacks-limits}
\externaldocument[stacks-cohomology-]{stacks-cohomology}
\externaldocument[stacks-perfect-]{stacks-perfect}
\externaldocument[stacks-introduction-]{stacks-introduction}
\externaldocument[stacks-more-morphisms-]{stacks-more-morphisms}
\externaldocument[stacks-geometry-]{stacks-geometry}
\externaldocument[moduli-]{moduli}
\externaldocument[moduli-curves-]{moduli-curves}
\externaldocument[examples-]{examples}
\externaldocument[exercises-]{exercises}
\externaldocument[guide-]{guide}
\externaldocument[desirables-]{desirables}
\externaldocument[coding-]{coding}
\externaldocument[obsolete-]{obsolete}
\externaldocument[fdl-]{fdl}
\externaldocument[index-]{index}

% Theorem environments.
%
\theoremstyle{plain}
\newtheorem{theorem}[subsection]{Theorem}
\newtheorem{proposition}[subsection]{Proposition}
\newtheorem{lemma}[subsection]{Lemma}

\theoremstyle{definition}
\newtheorem{definition}[subsection]{Definition}
\newtheorem{example}[subsection]{Example}
\newtheorem{exercise}[subsection]{Exercise}
\newtheorem{situation}[subsection]{Situation}

\theoremstyle{remark}
\newtheorem{remark}[subsection]{Remark}
\newtheorem{remarks}[subsection]{Remarks}

\numberwithin{equation}{subsection}

% Macros
%
\def\lim{\mathop{\mathrm{lim}}\nolimits}
\def\colim{\mathop{\mathrm{colim}}\nolimits}
\def\Spec{\mathop{\mathrm{Spec}}}
\def\Hom{\mathop{\mathrm{Hom}}\nolimits}
\def\Ext{\mathop{\mathrm{Ext}}\nolimits}
\def\SheafHom{\mathop{\mathcal{H}\!\mathit{om}}\nolimits}
\def\SheafExt{\mathop{\mathcal{E}\!\mathit{xt}}\nolimits}
\def\Sch{\mathit{Sch}}
\def\Mor{\mathop{\mathrm{Mor}}\nolimits}
\def\Ob{\mathop{\mathrm{Ob}}\nolimits}
\def\Sh{\mathop{\mathit{Sh}}\nolimits}
\def\NL{\mathop{N\!L}\nolimits}
\def\CH{\mathop{\mathrm{CH}}\nolimits}
\def\proetale{{pro\text{-}\acute{e}tale}}
\def\etale{{\acute{e}tale}}
\def\QCoh{\mathit{QCoh}}
\def\Ker{\mathop{\mathrm{Ker}}}
\def\Im{\mathop{\mathrm{Im}}}
\def\Coker{\mathop{\mathrm{Coker}}}
\def\Coim{\mathop{\mathrm{Coim}}}

% Boxtimes
%
\DeclareMathSymbol{\boxtimes}{\mathbin}{AMSa}{"02}

%
% Macros for moduli stacks/spaces
%
\def\QCohstack{\mathcal{QC}\!\mathit{oh}}
\def\Cohstack{\mathcal{C}\!\mathit{oh}}
\def\Spacesstack{\mathcal{S}\!\mathit{paces}}
\def\Quotfunctor{\mathrm{Quot}}
\def\Hilbfunctor{\mathrm{Hilb}}
\def\Curvesstack{\mathcal{C}\!\mathit{urves}}
\def\Polarizedstack{\mathcal{P}\!\mathit{olarized}}
\def\Complexesstack{\mathcal{C}\!\mathit{omplexes}}
% \Pic is the operator that assigns to X its picard group, usage \Pic(X)
% \Picardstack_{X/B} denotes the Picard stack of X over B
% \Picardfunctor_{X/B} denotes the Picard functor of X over B
\def\Pic{\mathop{\mathrm{Pic}}\nolimits}
\def\Picardstack{\mathcal{P}\!\mathit{ic}}
\def\Picardfunctor{\mathrm{Pic}}
\def\Deformationcategory{\mathcal{D}\!\mathit{ef}}
\usepackage{cleveref}
\crefname{originalreference}{original reference}{original references}
\Crefname{originalreference}{Original reference}{Original references}
\makeatletter
\newlabel{lemma-abelian-postnikov}{{0D9I}{source}{}{sourcepointer-1}{}}
\newlabel{lemma-abelian-postnikov@cref}{{[originalreference][0][]0D9I}{[1][0][]source}{}{}{}}
\makeatother
\crefname{originalreference}{original reference}{original references}
\Crefname{originalreference}{Original reference}{Original references}
\makeatletter
\newlabel{lemma-derived-cartesian-shriek}{{0D9E}{source}{}{sourcepointer-2}{}}
\newlabel{lemma-derived-cartesian-shriek@cref}{{[originalreference][0][]0D9E}{[1][0][]source}{}{}{}}
\makeatother
\crefname{originalreference}{original reference}{original references}
\Crefname{originalreference}{Original reference}{Original references}
\makeatletter
\newlabel{situation-simplicial-site}{{09WE}{source}{}{sourcepointer-3}{}}
\newlabel{situation-simplicial-site@cref}{{[originalreference][0][]09WE}{[1][0][]source}{}{}{}}
\makeatother
\title{Maps of simplicial systems lift to morphisms of cartesian objects under a vanishing assumption}
\author{The Stacks Project Authors}
\date{stacks-project:spaces-simplicial.tex:local-snapshot-20261012; GNU FDL 1.2 or later}
\begin{document}
\maketitle
\noindent\textit{Editorial note.} All delivered mathematical text is the author's own, copied from the pinned Stacks Project chapter file sources/stacks-spaces-simplicial/pinned-original-source.tex; no statement or proof step was rewritten, reordered or removed.\par\medskip
\noindent\textit{Editorial note.} The counted proof is the author's complete argument: the parenthetical '(omitted)' (source line 3122) covers the check that X\_1 -> X\_0 -> K' vanishes: X\_1 -> X\_0 is the alternating sum of the transitions and composing with the adjoint of K\_0 -> K'\_0 yields the alternating sum of the system maps, which vanishes by the definition of a morphism of simplicial systems - an unpacking of the adjunction; the closing 'this easily gives' (source line 3170) is the determination of the components gamma\_n from gamma\_0 through the transition isomorphisms of the cartesian object K, which forces \{gamma\_n\} to equal the given system morphism; both are auxiliary to the fully written induction.\par\medskip
\noindent\textit{Editorial note.} The label lemma-hom-cartesian-objects-derived names the counted statement, and the proof environment immediately following it in the source is delivered verbatim as its proof.\par\medskip
\noindent\textit{Editorial note.} The counted statement is carried by the author's own lemma environment in the chapter source.\par\medskip
\noindent\textit{Allowed standard prerequisites (editorial).}\begin{itemize}
\item The Postnikov system associated to a cartesian object. Same chapter, \textbackslash{}label\{lemma-abelian-postnikov\}. Lines 3115-3118: the objects X\_n and Postnikov system Y\_n attached to K; Y\_0 = X\_0 and the triangles Y\_1 -> X\_1 -> Y\_0 -> Y\_1[1].
\item Cartesian objects identified by the shriek-tensor maps. Same chapter, \textbackslash{}label\{lemma-derived-cartesian-shriek\} (card L46169 of this batch). Lines 3118-3120: cartesianness of K gives X\_n = g\_\{n!\}K\_n.
\item Adjunction between g\_\{n!\} and restriction to components. Same chapter, \textbackslash{}label\{lemma-restriction-to-components-site\} (card L46149). Lines 3133-3137: Hom(X\_\{n+1\}[n], K') = Hom(K\_\{n+1\}[n], K'\_\{n+1\}).
\item Morphisms of simplicial systems of the derived category. Same chapter, the definition of morphisms of simplicial systems. Lines 3115-3117 and 3120-3126: compatibility of the maps \{K\_n -> K'\_n\} with the system structure underlies the vanishing of X\_1 -> X\_0 -> K'.
\end{itemize}\medskip
\ifdefined\hypertarget\hypertarget{sourcepointer-1}{}\fi
\noindent\textit{Source reference.} Original reference 0D9I: sources/stacks-spaces-simplicial/pinned-original-source.tex, source lines 3002--3002. Prerequisite pointer invoked at proof line(s) 3119; it names a standard fact of the same chapter that the counted argument uses, and every step of the counted proof built on it is retained verbatim in the delivered body.\par\medskip
\ifdefined\hypertarget\hypertarget{sourcepointer-2}{}\fi
\noindent\textit{Source reference.} Original reference 0D9E: sources/stacks-spaces-simplicial/pinned-original-source.tex, source lines 2752--2752. Prerequisite pointer invoked at proof line(s) 3120; it names a standard fact of the same chapter that the counted argument uses, and every step of the counted proof built on it is retained verbatim in the delivered body.\par\medskip
\ifdefined\hypertarget\hypertarget{sourcepointer-3}{}\fi
\noindent\textit{Source reference.} Original reference 09WE: sources/stacks-spaces-simplicial/pinned-original-source.tex, source lines 888--888. Prerequisite pointer invoked at proof line(s) 3104; it names a standard fact of the same chapter that the counted argument uses, and every step of the counted proof built on it is retained verbatim in the delivered body.\par\medskip
\begin{lemma}
\label{lemma-hom-cartesian-objects-derived}
In Situation \ref{situation-simplicial-site}.
If $K, K' \in D(\mathcal{C}_{total})$.
Assume
\begin{enumerate}
\item $K$ is cartesian,
\item $\Hom(K_i[i - 1], K'_i) = 0$ for $i > 1$.
\end{enumerate}
Then any map $\{K_n \to K'_n\}$ between the associated simplicial systems 
of $K$ and $K'$ comes from a map $K \to K'$ in $D(\mathcal{C}_{total})$.
\end{lemma}

\begin{proof}
Let $\{K_n \to K'_n\}_{n \geq 0}$
be a morphism of simplicial systems of the derived category.
Consider the objects $X_n$ and Postnikov system $Y_n$
associated to $K$ of Lemma \ref{lemma-abelian-postnikov}.
By (1) and Lemma \ref{lemma-derived-cartesian-shriek} we have
$X_n = g_{n!}K_n$. In particular, the map $K_0 \to K'_0$
induces a morphism $X_0 \to K'$. Since $\{K_n \to K'_n\}$
is a morphism of systems, a computation (omitted) shows that
the composition
$$
X_1 \to X_0 \to K'
$$
is zero. As $Y_0 = X_0$ and as $Y_1$ fits into a distinguished
triangle
$$
Y_1 \to X_1 \to Y_0 \to Y_1[1]
$$
we conclude that there exists a morphism $Y_1[1] \to K'$ whose
composition with $X_0 = Y_0 \to Y_1[1]$ is the morphism $X_0 \to K'$
given above. Suppose given a map $Y_n[n] \to K'$ for $n \geq 1$.
From the distinguished triangle
$$
X_{n + 1}[n] \to Y_n[n] \to Y_{n + 1}[n + 1] \to X_{n + 1}[n + 1]
$$
we get an exact sequence
$$
\Hom(Y_{n + 1}[n + 1], K') \to
\Hom(Y_n[n], K') \to
\Hom(X_{n + 1}[n], K')
$$
As $X_{n + 1}[n] = g_{n + 1!}K_{n + 1}[n]$ the last group is equal to
$$
\Hom(K_{n + 1}[n], K'_{n + 1})
$$
which is zero by assumption (2). By induction we get a system of
maps $Y_n[n] \to K'$ compatible with transition maps and reducing
to the given map on $Y_0$. This produces a map
$$
\gamma :
K = \text{hocolim} Y_n[n]
\longrightarrow
K'
$$
This map in any case has the property that the diagram
$$
\xymatrix{
X_0 \ar[rd] \ar[r] &
K \ar[d]^\gamma \\
& K'
}
$$
is commutative. Restricting to
$\mathcal{C}_0$ we deduce that the map $\gamma_0 : K_0 \to K'_0$
is the same as the first map $K_0 \to K'_0$ of the morphism
of simplicial systems. Since $K$ is cartesian, this easily gives that
$\{\gamma_n\}$ is the map of simplicial systems we started out with.
\end{proof}

\end{document}
